\section{Related Work}\label{sec:related}

% TODO: this section is a starting point; extend before submission.

\paragraph{Non-interference} Non-interference goes back to Goguen and
Meseguer~\cite{goguen1982}; Sabelfeld and Myers survey language-based
approaches~\cite{sabelfeld2003}. Timing channels in programs have been closed by
padding transformations that equalise the running time of secret-dependent
branches~\cite{agat2000}; the sliced policy applies the same idea to handler
runs, with the NOP handler as padding.

\paragraph{Composable timing-sensitive non-interference} Our proofs rest on the
calculus of Rafnsson et al.~\cite{rafnsson2017}, whose definitions are
timing-sensitive by construction and whose theorems compose. We mechanise the
calculus and its composition theorems, and apply them to a model of interrupt
handling.

\paragraph{Verified operating systems} Machine-checked information-flow proofs
exist for real kernels, notably seL4~\cite{murray2013}, whose proof explicitly
leaves timing channels out of scope. Time protection~\cite{ge2019} argues that an
operating system must partition time as well as space, and includes the
partitioning of interrupts among its mechanisms\todo{check this characterisation
against~\cite{ge2019}}; our sliced policy is an instance of
that idea in a small model. What we add is a timing-sensitive, compositional,
mechanised proof that the mechanism is secure, for any non-interfering scheduler and
user processes.
